\documentclass[a4paper, 12pt]{article}

\usepackage[italian]{babel}
\usepackage{tikz}
\usepackage{xcolor}
\usepackage{graphicx}
\usepackage{hyperref}
\usepackage{imakeidx}
\usepackage{caption}
\usepackage{fancyhdr}
\usepackage{tabularx}

%--------------------VARIABILI-------------------- 
\def\logo{Immagini/logo.jpeg}
\def\ultima-versione{v1.0}
\def\titolo{Verbale interno}
\def\data{27/10/2024}
%------------------------------------------------

\usetikzlibrary{calc}
\definecolor{fp-blue}{HTML}{2885c8}
\definecolor{fp-red}{HTML}{ea5f64}
\makeindex[title=Indice]
\hypersetup{hidelinks}

\pagestyle{fancy}
\fancyhead[L]{}
\setlength{\headheight}{15pt}
\fancyhead[R]{\titolo - \ultima-versione}

\renewcommand{\familydefault}{\sfdefault}
\newcommand{\glossario}[1]{\fontfamily{lmr}\selectfont{\textit{#1\textsubscript{\small G}}}}

%--------------------INFORMAZIONI PRIMA PAGINA-------------------- 
\title{\Huge \textbf{\titolo}}
\author{\Large{Alt} \raisebox{0.3ex}{\normalsize  +} \Large{F4}}
\date{\data}
%----------------------------------------------------------------

\begin{document}

\begin{titlepage}      
    \maketitle
    \thispagestyle{empty}  

    \begin{tikzpicture}[remember picture, overlay]
        \fill[fp-blue] 
        ($(current page.south west) + (0, 10)$) 
        -- ($(current page.center) + (0, -8)$)
        -- ($(current page.center) + (0, -15)$)
        -- (current page.south west);

        \fill[fp-red]
        ($(current page.south east) + (0, 10)$) 
        -- ($(current page.center) + (0, -8)$)
        -- ($(current page.center) + (0, -15)$)
        -- (current page.south east);

        \clip ($(current page.center) + (0, -8)$) circle (1cm) node 
        {\includegraphics[width=.25\textwidth]{\logo}};
        
    \end{tikzpicture}    
\end{titlepage}

\tableofcontents

\newpage

\begin{table}[!h]
    \centering
    \caption*{\textbf{\Large Registro Modifiche}}
    {\renewcommand{\arraystretch}{2}
    \begin{tabularx}{\textwidth}{| X | X | X | X |}
        \hline
            \textbf{\large Versione} & 
            \textbf{\large Data} & 
            \textbf{\large Autore/i} & 
            \textbf{\large Descrizione} \\ 
        \hline \hline
            0.1 & 
            \data & 
            Eghosa Igbinedion & 
            Prima stesura del documento \\ 
        \hline 
            0.2 & 
            \data & 
            Marko Peric & 
            Revisione e modifica documento \\
        \hline 
    \end{tabularx}}
\end{table}

\newpage

\section{Registro presenze}
   \begin{itemize}
        \item[] \textbf{Ora inizio}: 18:00
        \item[] \textbf{Ora fine}: 19:30
        \item[] \textbf{Piattaforma}: Discord	
    \end{itemize}
\begin{table}[!h]
    \centering
    {\renewcommand{\arraystretch}{2}
    \begin{tabularx}{\textwidth}{| X | X |}
        \hline
            \textbf{\large Componente} & 
            \textbf{\large Presenza} \\ 
        \hline 
            Eghosa Igbinedion &
            Presente \\ 
        \hline 
            Giovanni Lan &
            Presente \\ 
        \hline 
            Enrico Bianchi &
            Presente \\ 
        \hline 
            Francesco Savio &
            Presente \\ 
        \hline 
            Marko Peric &
            Presente \\ 
        \hline 
            Pedro Leoni &
            Presente \\ 
        \hline 
    \end{tabularx}}
\end{table}

\newpage

\section{Way of working}
Il gruppo ha discusso sul workflow da adottare per il completamento dei requisiti minimi necessari alla candidatura a uno dei capitolati.

\subsection{Template}
Grazie alla maggiore esperienza di Pedro Leoni con lo strumento LaTeX, che si è deciso di usare per la stesura dei documenti, quest'ultimo si è offerto di creare un documento di esempio da utilizzare come template per i successivi verbali del gruppo.

\subsection{Gitflow}
Si è deciso di utilizzare gitflow come struttura di sviluppo per garantire un flusso di lavoro ordinato e scalabile, particolarmente utile in un progetto di team. Questo approccio permette di ridurre i conflitti e mantenere un elevato controllo sul codice in ogni fase del processo di sviluppo. La struttura si basa su branch principali:

\subsubsection{Main}
Contiene solo il codice stabile e pronto per il rilascio, assicurando che le versioni finali siano sempre affidabili.

\subsubsection{Develop}
Dove le funzionalità approvate dai branch di feature vengono integrate e testate, permettendo uno sviluppo continuo senza compromettere la stabilità del codice.

\subsubsection{Feature}
Creato per ogni nuova funzionalità, consente agli sviluppatori di lavorare in modo indipendente senza interferire con altre parti del progetto. Una volta completata, la feature viene unita a develop, garantendo così che solo le modifiche pronte siano integrate.

\subsection{Organizzazione della repository}
Si è concordato di sfruttare il sistema di \textbf{Issues}, il sistema di \textbf{Milestone} e il sistema di \textbf{Tables} per:
\begin{itemize}
    \item Tracciare e assegnare compiti e responsabilità.
    \item Monitorare lo stato di avanzamento dei requisiti.
    \item Creare compiti non assegnati che possono essere svolti in caso di completamento anticipato dei propri compiti.
\end{itemize}
Ogni membro del gruppo sarà incaricato di aprire e gestire le issue legate ai propri compiti, garantendo trasparenza e tracciabilità dei progressi.

\section{Brainstorming per le riunioni esterne}
Durante la riunione, il gruppo si è focalizzato sulla preparazione di domande mirate a comprendere in maniera più approfondita la natura dei capitolati C5, C7 e C9, precedentemente selezionati come i più rilevanti. Le domande sono state suddivise in quesiti di carattere generale, comuni a tutte le aziende, e domande tecniche specifiche per ogni capitolato. Di seguito sono riportati alcuni esempi.

\subsection{Domande generali}
\begin{itemize}
    \item Quali canali di comunicazione si intendano utilizzare per mantenere il contatto tra azienda e gruppo di sviluppo;
    \item La frequenza di interazione prevista tra azienda e gruppo durante le diverse fasi del progetto;
    \item La piattaforma di destinazione per il rilascio e la fruizione del prodotto finale;
    \item Quali componenti o aspetti del progetto, secondo il parere degli sviluppatori aziendali, potrebbero presentare le maggiori difficoltà tecniche.
\end{itemize}

\subsection{Domande tecniche}
\subsubsection{Azzurro Digitale (C9)}
\begin{itemize}
    \item Chiarimenti in merito alla menzione della "documentazione di API di terze parti" nel capitolato;
    \item Informazioni sulle risorse eventualmente fornite dall'azienda per l'accesso alle API di terze parti;
    \item Preferenze specifiche riguardo all'uso di Node.js o Spring Boot per lo sviluppo del progetto;
    \item Precisazioni sulla funzione principale della libreria \textit{Langchain} e sul suo ruolo previsto nel progetto.
\end{itemize}

\subsubsection{Ergon Informatica (C7)}
\begin{itemize}
    \item Chiarimenti riguardo alla necessità di sviluppare un ambiente per la gestione di Large Language Model (LLM) o se tale ambiente venga fornito dall'azienda;
    \item Dettagli sulla gestione e struttura dei Database Vettoriali richiesti dal capitolato.
\end{itemize}

\subsubsection{SanMarco Informatica (C5)}
\begin{itemize}
    \item Se per l’apprendimento delle librerie \textit{three.js} e \textit{d3.js} siano previste risorse aziendali o se tale apprendimento debba avvenire in modo autonomo; \item Eventuali restrizioni o linee guida specifiche per lo sviluppo del \textit{backend} del progetto. \end{itemize}

\section{Conclusioni}
L’incontro si conclude con la conferma della necessità di ulteriori approfondimenti sui capitolati scelti e l’organizzazione di una prossima riunione per valutare le risposte ricevute dalle aziende.
\end{document}